\practicechapter{第 0 章实践：学习地图与证据链}

本章对应第一册第 0 章“学习地图：把程序、机器、系统和性能连成一条主线”。正文讲的是心智模型：不要把 C++、编译器、CPU、操作系统和性能测量分成互不相干的课。本章要把这个心智模型落成实践卷的总地图。最终你要能回答一个朴素问题：给定一个自己编译出的可执行文件，哪些结论来自文件字节，哪些结论来自工具解释，哪些结论来自运行时观察，哪些结论还只是猜测。

\practicesection{一、对应原书问题：主线不是名词表}

如果只说“源码经过编译变成机器码，然后 CPU 执行”，读者很容易以为学习路线就是背名词。实践卷从一开始就要求把名词绑定到证据。源码文件是输入，编译命令是转换规则，ELF 文件是二进制产物，section 和 symbol 是文件里的结构化证据，loader 与进程是运行时事实，benchmark 是在特定环境下观察到的数字。

本卷贯穿项目 \term{Executable Evidence Kit}{可执行文件证据工具箱} 的第一份交付不是代码，而是一张证据链表。表中至少要有五列：对象、能观察到的字段、观察命令、代码中的对应结构、当前限制。例如 \code{ELF header} 的字段可以由 \code{cpu1ee\_cli} 和 \cmd{readelf -h} 同时观察；\code{.text} 的 offset 和 size 可以由 section table 给出；真正装载后的虚拟地址则不能只靠文件 offset 推断。

\begin{keyidea}
实践卷按原书章节走，但所有章节都围绕同一个工程对象。这样做的目的不是重复讲一遍概念，而是让每个概念都能落在同一份二进制证据上。
\end{keyidea}

\practicesection{二、从特例推到规律：文件名不是证据}

先看一个特例。你编译出 \code{hello}，随后重新编译一次，文件名仍然叫 \code{hello}。如果报告只写“我分析了 hello”，这份报告就没有固定分析对象，因为第二次构建已经覆盖了第一次构建。再看另一个特例：同一个源码用 Debug 和 Release 构建，文件名也许只差一个目录，但 section size、symbol 数量、反汇编和运行时间都可能不同。

由这两个特例可以推出第一条规律：实践报告必须绑定字节身份，而不是绑定口头描述。本卷用 \code{read\_binary\_file} 读入字节，用 \code{checksum\_bytes} 计算摘要，用 byte size 记录长度。checksum 不是密码学证明，它的作用是把“这一次输入”固定下来，避免报告描述错对象。

第二条规律是区分文件事实和运行事实。ELF header、section table、symbol table 属于文件事实；进程地址空间、动态链接器状态、perf 事件属于运行事实。文件事实可以离线读取，运行事实必须在程序执行时观察。把这两类混在一起，就会把 section offset 误当成进程地址，或者把一次 benchmark 数字误当成二进制结构结论。

\practicesection{三、实践任务：建立项目总地图}

先完成一次完整构建，确认后续章节使用的是同一份工程：

\begin{lstlisting}[language=bash]
cmake -S books/cpu-volume-1-practice \
      -B books/cpu-volume-1-practice/build/executable-evidence-debug \
      -DCMAKE_BUILD_TYPE=Debug
cmake --build books/cpu-volume-1-practice/build/executable-evidence-debug
ctest --test-dir books/cpu-volume-1-practice/build/executable-evidence-debug \
      --output-on-failure
\end{lstlisting}

然后阅读三个入口文件：\filepath{include/cpu1/executable\_evidence.hpp}、\filepath{src/executable\_evidence.cpp}、\filepath{src/main.cpp}。不要急着改代码，先写一张表，把结构体和概念对应起来：\code{ElfReport} 对应一次检查报告，\code{SectionSummary} 对应 section 摘要，\code{SymbolSummary} 对应符号摘要，\code{InspectionConfig} 对应解析边界。

\begin{labbox}
本章交付物是 \filepath{reports/ch00-evidence-map.md}。报告中必须列出至少 8 个证据点：source name、byte size、checksum、ELF magic、machine、entry address、\code{.text} section、symbol、diagnostic、benchmark smoke 可以任选其八，但每一项都要写出它来自文件、工具、测试还是运行。
\end{labbox}

\practicesection{四、验收：能说明边界才算会用地图}

本章验收不看你画得漂不漂亮，只看证据边界是否清楚。通过标准如下：

\begin{enumerate}
  \item 能运行 \cmd{ctest}，并说明测试通过只证明当前合同，没有证明所有 ELF 都支持。
  \item 能指出 \code{checksum} 固定的是输入字节，不是源码语义，也不是程序运行行为。
  \item 能说明 \code{entry address} 不是 \code{main} 的地址。
  \item 能区分 section 的 file offset 和进程中的 virtual address。
  \item 能写出本卷第一版不支持的内容：32 位 ELF、大端 ELF、PE/Mach-O、DWARF、动态链接装载后地址。
\end{enumerate}

如果这些边界讲不清，后续章节很容易变成“看见一个字段就下结论”。第 0 章实践的目的就是把这种风险提前挡住。
